% ============================================================
%  MTH636 -- Game Theory | Lecture 8
%  Iterated Elimination and Nash Equilibrium
% ============================================================
\documentclass[12pt, a4paper]{article}

% ============================================================
%  MTH636 -- Game Theory
%  Notes Preamble — Color-Coded, TikZ-Enhanced
% ============================================================

% --- Core packages ---
\usepackage[T1]{fontenc}
\usepackage[utf8]{inputenc}
\usepackage{lmodern}
\usepackage[margin=2.5cm, top=2.8cm, bottom=2.8cm]{geometry}
\usepackage{microtype}

% --- Math ---
\usepackage{amsmath, amssymb, amsthm, mathtools}
\usepackage{bm}        % bold math
\usepackage{cancel}    % strike-through in math

% --- Color & Boxes ---
\usepackage[dvipsnames, svgnames, x11names, table]{xcolor}
\usepackage[most]{tcolorbox}
\tcbuselibrary{breakable, skins, theorems}

% --- TikZ & Diagrams ---
\usepackage{tikz}
\usepackage{tikz-qtree}
\usetikzlibrary{arrows.meta, positioning, shapes.geometric, fit, backgrounds,
                decorations.pathreplacing, calc, matrix, trees}
\usepackage{pgfplots}
\pgfplotsset{compat=1.18}

% --- Tables ---
\usepackage{booktabs}
\usepackage{multirow}
\usepackage{array}
\usepackage{tabularx}
\usepackage{makecell}

% --- Typography & Layout ---
\usepackage{enumitem}
\usepackage{parskip}
\usepackage{fancyhdr}
\usepackage{titlesec}
\usepackage{titletoc}
\usepackage{hyperref}
\usepackage{graphicx}
\usepackage{float}
\usepackage{caption}

% --- Color Palette ---
\definecolor{defblue}{HTML}{1A5276}      % Definition header
\definecolor{defbluebg}{HTML}{D6EAF8}    % Definition background
\definecolor{thmgreen}{HTML}{1E8449}     % Theorem header
\definecolor{thmbg}{HTML}{D5F5E3}        % Theorem background
\definecolor{exorange}{HTML}{B7770D}     % Example header
\definecolor{exbg}{HTML}{FEF9E7}         % Example background
\definecolor{warnred}{HTML}{922B21}      % Warning/Key Insight header
\definecolor{warnbg}{HTML}{FADBD8}       % Warning background
\definecolor{noticegray}{HTML}{2C3E50}   % Note header
\definecolor{noticebg}{HTML}{EBF5FB}     % Note background
\definecolor{taketawayyellow}{HTML}{7D6608} % Key Takeaway header
\definecolor{takeawaybg}{HTML}{FDFEFE}   % Takeaway background
\definecolor{accentpurple}{HTML}{6C3483} % Accent / proof
\definecolor{lightgray}{HTML}{F2F3F4}    % shading

% ============================================================
%  TCOLORBOX ENVIRONMENTS
% ============================================================

% -- Definition (Blue) --
\newtcolorbox{defbox}[2][]{
  enhanced, breakable,
  colback=defbluebg, colframe=defblue, coltitle=white,
  fonttitle=\bfseries\sffamily,
  title={Definition: #2},
  attach boxed title to top left={yshift=-2mm, xshift=4mm},
  boxed title style={colback=defblue, rounded corners},
  arc=4pt, boxrule=1.2pt,
  left=6pt, right=6pt, top=8pt, bottom=6pt,
  #1
}

% -- Theorem / Proposition (Green) --
\newtcolorbox{thmbox}[2][]{
  enhanced, breakable,
  colback=thmbg, colframe=thmgreen, coltitle=white,
  fonttitle=\bfseries\sffamily,
  title={Theorem: #2},
  attach boxed title to top left={yshift=-2mm, xshift=4mm},
  boxed title style={colback=thmgreen, rounded corners},
  arc=4pt, boxrule=1.2pt,
  left=6pt, right=6pt, top=8pt, bottom=6pt,
  #1
}

% -- Proposition (Green variant) --
\newtcolorbox{propbox}[2][]{
  enhanced, breakable,
  colback=thmbg, colframe=thmgreen!70!black, coltitle=white,
  fonttitle=\bfseries\sffamily,
  title={Proposition: #2},
  attach boxed title to top left={yshift=-2mm, xshift=4mm},
  boxed title style={colback=thmgreen!70!black, rounded corners},
  arc=4pt, boxrule=1.2pt,
  left=6pt, right=6pt, top=8pt, bottom=6pt,
  #1
}

% -- Example (Orange/Yellow) --
\newtcolorbox{exbox}[2][]{
  enhanced, breakable,
  colback=exbg, colframe=exorange, coltitle=white,
  fonttitle=\bfseries\sffamily,
  title={Example: #2},
  attach boxed title to top left={yshift=-2mm, xshift=4mm},
  boxed title style={colback=exorange, rounded corners},
  arc=4pt, boxrule=1.2pt,
  left=6pt, right=6pt, top=8pt, bottom=6pt,
  #1
}

% -- Key Insight / Warning (Red) --
\newtcolorbox{warnbox}[2][]{
  enhanced, breakable,
  colback=warnbg, colframe=warnred, coltitle=white,
  fonttitle=\bfseries\sffamily,
  title={Key Insight: #2},
  attach boxed title to top left={yshift=-2mm, xshift=4mm},
  boxed title style={colback=warnred, rounded corners},
  arc=4pt, boxrule=1.2pt,
  left=6pt, right=6pt, top=8pt, bottom=6pt,
  #1
}

% -- Note (Gray-Blue) --
\newtcolorbox{notebox}[1][]{
  enhanced, breakable,
  colback=noticebg, colframe=noticegray!60, coltitle=white,
  fonttitle=\bfseries\sffamily,
  title={Note},
  attach boxed title to top left={yshift=-2mm, xshift=4mm},
  boxed title style={colback=noticegray!60, rounded corners},
  arc=4pt, boxrule=0.8pt,
  left=6pt, right=6pt, top=8pt, bottom=6pt,
  #1
}

% -- Key Takeaway --
\newtcolorbox{takebox}[1][]{
  enhanced, breakable,
  colback=takeawaybg, colframe=taketawayyellow, coltitle=white,
  fonttitle=\bfseries\sffamily,
  title={Key Takeaway},
  attach boxed title to top left={yshift=-2mm, xshift=4mm},
  boxed title style={colback=taketawayyellow, rounded corners},
  arc=4pt, boxrule=1.2pt,
  left=6pt, right=6pt, top=8pt, bottom=6pt,
  #1
}

% -- Intuition --
\newtcolorbox{intubox}[1][]{
  enhanced, breakable,
  colback=accentpurple!8, colframe=accentpurple!60, coltitle=white,
  fonttitle=\bfseries\sffamily,
  title={Intuition},
  attach boxed title to top left={yshift=-2mm, xshift=4mm},
  boxed title style={colback=accentpurple!80, rounded corners},
  arc=4pt, boxrule=0.8pt,
  left=6pt, right=6pt, top=8pt, bottom=6pt,
  #1
}

% ============================================================
%  PAYOFF MATRIX HELPER
%  Usage: inside tabular, use \payoff{a}{b} for (a,b)
% ============================================================
\newcommand{\payoff}[2]{$(#1,\;#2)$}
\newcommand{\nash}[2]{\underline{$#1$},\underline{$#2$}}  % underline NE payoffs
\newcommand{\Nset}{\mathbb{N}}
\newcommand{\Rset}{\mathbb{R}}
\newcommand{\Sset}{\mathcal{S}}
\newcommand{\Ui}{u_i}

% ============================================================
%  SECTION FORMATTING
% ============================================================
\titleformat{\section}
  {\large\bfseries\sffamily\color{defblue}}
  {\thesection}{1em}{}[\titlerule]

\titleformat{\subsection}
  {\normalsize\bfseries\sffamily\color{thmgreen!80!black}}
  {\thesubsection}{1em}{}

\titleformat{\subsubsection}
  {\normalsize\itshape\sffamily\color{exorange!80!black}}
  {\thesubsubsection}{1em}{}

% ============================================================
%  HEADER / FOOTER
% ============================================================
\pagestyle{fancy}
\fancyhf{}
\fancyhead[L]{\sffamily\small\color{defblue} MTH636 -- Game Theory}
\fancyhead[R]{\sffamily\small\color{defblue} \leftmark}
\fancyfoot[C]{\sffamily\small\thepage}
\renewcommand{\headrulewidth}{0.4pt}

% ============================================================
%  HYPERREF SETUP
% ============================================================
\hypersetup{
  colorlinks=true,
  linkcolor=defblue,
  urlcolor=thmgreen,
  citecolor=accentpurple,
  bookmarksnumbered=true,
  pdftitle={MTH636 Game Theory Notes},
  pdfauthor={MTH636 Student Notes}
}

% ============================================================
%  MISC HELPERS
% ============================================================
\newcommand{\BR}{\text{BR}}        % Best Response
\newcommand{\NE}{\text{NE}}        % Nash Equilibrium
\newcommand{\SPNE}{\text{SPNE}}    % Subgame Perfect NE
\newcommand{\Dom}{\text{dom}}      % Dominance
\newcommand{\mixs}[1]{\sigma_{#1}} % Mixed strategy
\newcommand{\pures}[1]{s_{#1}}     % Pure strategy


% ============================================================
\begin{document}

% Title
\begin{center}
  {\LARGE\bfseries\sffamily\color{defblue} MTH636 -- Game Theory}\\[6pt]
  {\Large\sffamily Lecture 8: Iterated Elimination \& Nash Equilibria}\\[4pt]
\end{center}
\vspace{0.5cm}
\hrule height 1.5pt
\vspace{0.3cm}

% ============================================================
\section{Iterated Elimination of Strictly Dominated Strategies (IESDS)}

Recall that for a game $G = \langle N, (S_i)_{i \in N}, (u_i)_{i \in N} \rangle$, a strategy profile $s^*$ is a Nash Equilibrium if $u_i(s^*) \geq u_i(s_i, s^*_{-i})$ for all $i \in N, s_i \in S_i$.

Suppose we create a new game $\hat{G} = \langle N, (S'_i)_{i \in N}, (u_i)_{i \in N} \rangle$ by removing a \textbf{strictly dominated strategy} $\bar{s}_k$ for some player $k$. So:
\[
  S_i' = S_i \text{ for } i \neq k, \quad \text{and} \quad S_k' = S_k \setminus \{\bar{s}_k\}
\]

\begin{thmbox}{IESDS and Nash Equilibria}
  The set of Nash Equilibria in $G$ and in $\hat{G}$ remain \textbf{exactly the same}.
  \begin{enumerate}
      \item \textbf{IESDS cannot remove any Nash Equilibrium:} If an NE were removed, it must have had $\bar{s}_k$ as its $k$-th component. But a strictly dominated strategy can never be a best response, so it can never be part of a Nash Equilibrium.
      \item \textbf{IESDS cannot create new Nash Equilibria:} Could a profile $s^*$ be an NE in $\hat{G}$ but not in $G$? If it wasn't an NE in $G$, some player must have had a profitable deviation. Since all strategies except $\bar{s}_k$ are still in $\hat{G}$, the only way the deviation is missing is if the profitable deviation was to $\bar{s}_k$ itself! But $\bar{s}_k$ is strictly dominated by some $t_k \in S'_k$, meaning $t_k$ would be an even better deviation that is still available in $\hat{G}$. Thus, no new NE is created.
  \end{enumerate}
\end{thmbox}

% ============================================================
\section{Iterated Elimination of Weakly Dominated Strategies}

The strong guarantee of IESDS \textbf{does not hold} for weakly dominated strategies.

\begin{warnbox}{Weakly Dominated Strategies can remove NE}
  Consider the following game:
  \begin{center}
  \renewcommand{\arraystretch}{1.3}
  \begin{tabular}{c|cc}
    & L & R \\
    \hline
    T & \payoff{0}{0} & \payoff{2}{1} \\
    B & \payoff{3}{2} & \payoff{1}{2} \\
  \end{tabular}
  \end{center}
  \textbf{Nash Equilibria of Original Game:} Both $(B, L)$ and $(T, R)$ are Nash Equilibria.
  
  \textbf{Elimination:}
  \begin{itemize}
      \item For Player 2, R weakly dominates L (payoffs: $1 \geq 0$ and $2 \geq 2$).
      \item If we eliminate L, we are left with the right column.
      \item In the reduced game (T vs R yields 2, B vs R yields 1), T strictly dominates B for Player 1. We eliminate B.
      \item We are left with only $(T, R)$.
  \end{itemize}
  By eliminating weakly dominated strategies, we \textbf{lost} the Nash Equilibrium $(B, L)$!
\end{warnbox}

\textbf{Question:} Can iterated elimination of weakly dominated strategies create \textit{new} Nash Equilibria that were not NE in the original game? (No, but proving it requires care).

% ============================================================
\section{Introduction to Mixed Strategies}

Until now, we assumed players choose a specific strategy deterministically (Pure Strategies). But what if players can randomize?

Consider a choice between:
\begin{itemize}
    \item Option A: Receive Rs 100 for sure.
    \item Option B: Receive Rs 500 with probability $1/4$, and Rs 0 with probability $3/4$.
\end{itemize}
We assume players are \textbf{rational} and evaluate probability distributions using \textbf{Expected Utility}.
The expected value of B is $\frac{1}{4}(500) + \frac{3}{4}(0) = 125$. If utility is linear in money, the player prefers B over A.

\subsection{Affine Transformations of Utilities}
Suppose we have a game $G$ with payoffs $u_i$. We create an equivalent game $\hat{G}$ by applying a function to the utilities: $\hat{u}_i = v_i \circ u_i$.
\begin{itemize}
    \item If $v_i$ is a general strictly increasing monotonic transformation, it preserves preferences over \textbf{pure} strategies, but it \textbf{does not} preserve preferences over \textbf{mixed} (randomized) strategies, because it changes the expected values relative to each other.
    \item To preserve preferences over mixed strategies, we must only allow \textbf{Positive Affine Transformations}:
    \[
      v_i(x) = a_i x + b_i \quad \text{where } a_i > 0
    \]
\end{itemize}

\end{document}